%! Comparing Distributions — Unit 2, Lesson 2.4 — Warm-Up
% ONE PAGE, blank and key. Three items at 12pt (item spacing 22pt) so each item gets real
% handwriting room and still fits. Do not add a fourth item — the page is full.
% \workrowsep is 30pt here (not 1.2's 50pt): the page also carries two fill tables.
% GRADUAL RELEASE: the warm-up activates prior knowledge before the guided notes. Items 1 and 2
% are Lesson 2.3's quartiles, IQR and 1.5 IQR fence on TONY'S PIZZA — the same 11 delivery
% times the notes open with (row 1's back-to-back stemplot, row 2's boxplot with 44 plotted as
% an outlier). Item 3 is two tiny lists with the SAME mean and different spread: it seeds the
% crux (row 4, "Different Spread") without the vocabulary.
% THE DATA — Tony's: 18 21 22 24 25 25 27 28 30 33 44 (n=11). Median = 6th = 25; lower half
%   18 21 22 24 25 -> Q1 22; upper half 27 28 30 33 44 -> Q3 30; IQR 8; upper fence
%   30 + 12 = 42 -> 44 is an outlier. Lists: A 9 10 11 (mean 10, range 2); B 2 10 18
%   (mean 10, range 16).
\documentclass[12pt]{article}
\usepackage{saar-article}
\usepackage{saar-boxes}
\newcommand{\TallMath}[1]{$\displaystyle #1\rule[-1.4em]{0pt}{3.2em}$}
% Word bank strip for every fill-in cluster. Defined per document; the key inherits it and
% prints the same strip, so heights match.
\newcommand{\wordbank}[1]{\par\vspace{2pt}\noindent\fcolorbox{goldacc}{goldbg}{\parbox{\dimexpr\linewidth-2\fboxsep-2\fboxrule\relax}{\small\textbf{Word bank:} #1}}\par\vspace{4pt}}
% Handwriting room inside every \begin{work} block. Moves the blank and the
% key together, so the two can never drift apart.
\setlength{\workrowsep}{30pt}

\begin{document}

\pageheader{Unit 2, Lesson 2.4}{Warm-Up}
% NAMESTRIP: no \namedateperiod here. The student writes their name once, on the
% cover this component is stapled behind.

\vspace{0.15in}

\boxguard
\begin{notesbox}{Spiral Review --- Tony's Pizza Delivery Times}
\normalsize
Tony's Pizza timed its last $11$ deliveries, in minutes, in order:
\quad $18,\ 21,\ 22,\ 24,\ 25,\ 25,\ 27,\ 28,\ 30,\ 33,\ 44$.

\begin{enumerate}[leftmargin=1.8em, itemsep=22pt, topsep=10pt]

  \item Find the quartiles and the IQR. (Quartiles are the medians of each half; with $11$
        values, leave the overall median out of both halves.)

        \vspace{4pt}
        \renewcommand{\arraystretch}{2.4}
        \begin{tabularx}{\linewidth}{@{} Y Y Y Y @{}}
        \toprule
        $Q_1$ & \textbf{Median} & $Q_3$ & \textbf{IQR} \\
        \midrule
        \blank{2.2cm} & \blank{2.2cm} & \blank{2.2cm} & \blank{2.2cm} \\
        \bottomrule
        \end{tabularx}

  \item Is the $44$-minute delivery an outlier? Use the upper fence,
        $Q_3 + 1.5 \cdot \text{IQR}$.

        \begin{work}
          \text{Upper fence} &= 30 + 1.5(8) \\
                             &= 42 \text{ minutes}
        \end{work}

        \vspace{6pt}
        Outlier? \blank{10cm}

  \item Two tiny lists: \quad \textbf{List A:} $9,\ 10,\ 11$ \qquad \textbf{List B:} $2,\ 10,\ 18$.
        \wordbank{2 \;\textbullet\; 10 \;\textbullet\; 16 \;\textbullet\; 20}

        \vspace{4pt}
        \renewcommand{\arraystretch}{2.2}
        \begin{tabularx}{\linewidth}{@{} X p{3.4cm} p{3.4cm} @{}}
        \toprule
        \textbf{List} & \textbf{Mean} & \textbf{Range} \\
        \midrule
        A: $9,\ 10,\ 11$  & \blank{3cm} & \blank{3cm} \\
        B: $2,\ 10,\ 18$  & \blank{3cm} & \blank{3cm} \\
        \bottomrule
        \end{tabularx}

        \vspace{8pt}
        In one sentence: how are the two lists alike, and how are they different?
        \par\writelines{2}

\end{enumerate}
\end{notesbox}

\end{document}
